% Part: proof-theory
% Chapter: natural-deduction
% Section: grafting

\documentclass[../../../include/open-logic-section]{subfiles}

\begin{document}

\olfileid{pt}{nat}{gra}

\olsection{Grafting \usetoken{P}{proof}}

!!^{proof}s in \Log{N1ci} are !!{proof}s of !!{formula}s~$!A$ from
assumptions~$!B$. Suppose $!B$ itself has !!a{proof}. How do we
combine the !!{proof} of~$!B$ with the !!{proof} of $!A$ from~$!B$ to
get !!a{proof} of~$!A$ that doesn't use the assumption~$!B$?

One option is to remove the assumption~$!B$ from the !!{proof} of~$!A$
by applying~\Intro{\lif}. Then we can combine the resulting !!{proof}
of~$!B \lif !A$ with the !!{proof} of~$!B$ to obtain
\[
\AxiomC{$\Discharge{!B}{x}$}
\DeduceC{$!A$}
\DischargeRule{\Intro{\lif}}{x}
\UnaryInfC{$!B \lif !C$}
\AxiomC{}
\DeduceC{$!B$}
\RightLabel{\Elim{\lif}}
\BinaryInfC{$!A$}
\DisplayProof
\]
This is straightforward, but somewhat ugly: why introduce $\lif$ only
to eliminate it right away? It should be possible to avoid this kind
of ``detour'' through $!B \lif !A$. (That they can always be avoided
is in fact the content of the normalization theorem.)

In \Log{N1ci} we can actually avoid it in this case quite easily: we
can simply replace the assumptions~$!B^x$ by the entire !!{proof}
of~$!B$. This ``grafting'' of !!{proof}s onto the assumptions of
others is quite useful and so we record it as a definition.

\begin{defn}
If $\delta_1$ is a \Log{N1ci}-!!{proof} of~$!A$ from open
assumption~$!B^x$, and $\delta_2$ is a \Log{N1ci}-!!{proof} of~$!B$,
then $\Subst{\delta_1}{\delta_2}{!B^x}$ is the result of replacing
every undischarged occurrence of~$!B^x$ in~$\delta_1$ by~$\delta_2$.
\end{defn}

Provided that $\delta_1$ and $\delta_2$ don't have different
assumption !!{formula}s labelled with the same label, and no open
assumption of~$\delta_2$ contains an eigenvariable of~$\delta_1$, the
resulting $\Subst{\delta_1}{\delta_2}{!B^x}$ is a correct !!{proof},
and its open assumptions are those of $\delta_1$ together with those
of~$\delta_2$. When grafting !!{proof}s, these conditions will
typically be satisfied. The first is satisfied whenever $\delta_1$ and
$\delta_2$ are subproofs of a larger !!{proof}. If the second is not
satisfied, we can rename the eigenvariables of~$\delta_1$ so that it
is.

A corresponding definition applies to \Log{N2ci}.

\begin{defn}
If $\delta_1$ is a \Log{N2ci}-!!{proof} of~$x:!B, \Gamma_1 \Sequent
!A$, and $\delta_2$ is a \Log{N1ci}-!!{proof} of~$\Gamma_2 \Sequent
!B$, then $\Subst{\delta_1}{\delta_2}{x:!B}$ is the result of
replacing every axiom $x:!B \Sequent !B$ in~$\delta_1$ by~$\delta_2$,
and adding $\Gamma_2$ to the context of every sequent below such
axioms.
\end{defn}

\begin{prop}
If $\delta_1 \Proves x:!B, \Gamma_1 \Sequent
!A$ and $\delta_2 \Proves \Gamma_2 \Sequent !B$, then
$\Subst{\delta_1}{\delta_2}{x:!B} \Proves \Gamma_1, \Gamma_2 \Sequent
!A$, provided no eigenvariable of~$\delta_1$ occurs in~$\Gamma_2$.
\end{prop}


\begin{proof}
  By induction on $\pheight{\delta}$.
\end{proof}

\end{document}